% Machine-checked reflection positivity for finite abelian lattice gauge theory.
% Lane: O (operator bridge), phase O-3, modules YangMills/OS/**.
% REGIME: tricotomy; NO scores, NO desk language, NO commit chronology beyond
% the reproducibility freeze.
\documentclass[11pt]{article}
\usepackage[margin=1.1in]{geometry}
\usepackage{amsmath,amssymb,amsthm}
\usepackage{booktabs,longtable,array}
\usepackage[colorlinks=true,linkcolor=blue,urlcolor=blue,citecolor=blue]{hyperref}

\newtheorem{theorem}{Theorem}[section]
\newtheorem{lemma}[theorem]{Lemma}
\newtheorem{proposition}[theorem]{Proposition}
\newtheorem{corollary}[theorem]{Corollary}
\newtheorem{remark}[theorem]{Remark}
\theoremstyle{definition}
\newtheorem{definition}[theorem]{Definition}

\emergencystretch=3em
\tolerance=2000

\newcommand{\lean}[1]{\texttt{#1}}
\newcommand{\anchor}{55d346027a5181157aa1f5683b97247d115801f1}
\newcommand{\osfile}[2]{%
  \href{https://github.com/lluiseriksson/THE-ERIKSSON-PROGRAMME/blob/\anchor/YangMills/OS/#1.lean}{\nolinkurl{#2}}}
\newcommand{\repofile}[2]{%
  \href{https://github.com/lluiseriksson/THE-ERIKSSON-PROGRAMME/blob/\anchor/#1}{\nolinkurl{#2}}}
\newcommand{\R}{\mathbb{R}}
\newcommand{\C}{\mathbb{C}}
\newcommand{\Z}{\mathbb{Z}}
\newcommand{\SU}{\mathrm{SU}}
\newcommand{\inn}[2]{\langle #1,\, #2\rangle}

\title{Machine-Checked Reflection Positivity\\
for Finite Abelian Lattice Gauge Theory}
\author{Lluis Eriksson}
\date{July 27, 2026 (version 1.0)}

\begin{document}
\maketitle

\begin{abstract}
We machine-check, in Lean~4 with no \lean{sorry} and no project axioms, the
Osterwalder--Seiler reflection positivity of a lattice gauge theory with finite
abelian gauge group.  The development is organised so that the three
ingredients are separated and each is proved on its own: an analytic step, a
geometric step, and the single place where a property of the Boltzmann factor
is actually used.

The analytic step is that a crossing kernel of the form
$K(x,y) = \sum_i c_i\,\varphi_i(x)\overline{\varphi_i(y)}$ with $c_i \ge 0$ is
positive semidefinite, that this class is closed under products, and that it is
closed under conjugation by a positive diagonal.  Formulating the hypothesis as
a non-negative combination of characters rather than as non-negativity of
Fourier coefficients removes any need for Bochner's theorem on a finite abelian
group and for the Schur product theorem: the development uses no spectral and
no matrix-positivity API.

The geometric step is a splitting of the configuration space across the
reflection plane under which the reflection is the swap and the Gibbs weight
factors as $w(x)\,w(y)\,K(x,y)$.  We prove that the Osterwalder--Seiler pairing
of an observable of one half against its reflection is then \emph{exactly} the
quadratic form of $w(x)K(x,y)w(y)$, so that reflection positivity follows from
the analytic step.

The physical step is the instance.  For $\Z_2$ the Wilson factor
$e^{\beta s}$, $s = \pm 1$, expands in the two characters with coefficients
$(e^{\beta} \pm e^{-\beta})/2$, both non-negative exactly when $\beta \ge 0$;
so the $\Z_2$ Wilson crossing kernel is positive semidefinite at non-negative
coupling.  For $\Z_N$ with $N > 2$ the coefficients are discrete Bessel-type
sums and their non-negativity is \emph{not} established here.

We are explicit about what is absent: no Gelfand--Naimark--Segal quotient, no
transfer operator, no identification of a Euclidean correlator with a matrix
element, and therefore no mass gap.  Nothing here is a claim about $\SU(N)$,
the continuum limit, or the Clay problem.
\end{abstract}

\section{What is proved, and what is not}

Reflection positivity is the hypothesis from which Osterwalder and
Seiler~\cite{OsterwalderSeiler} construct the physical Hilbert space of a
lattice gauge theory and its transfer operator.  It is the first of four steps
between a Euclidean lattice measure and a statement about a spectrum, the
others being the Gelfand--Naimark--Segal quotient, the transfer operator
itself, and the identification of Euclidean correlators with its matrix
elements.

This paper mechanizes the first step for a finite abelian gauge group, and
\emph{only} the first step.  In particular the following are not established
here, nor anywhere in the Lean development accompanying this paper: the GNS
quotient; the existence of a transfer operator for any gauge theory; the
identification $\mathbb{E}[A\cdot\theta_n B] = \inn{A\Omega}{T^n B\Omega}$; and
any statement whatever about a mass gap.  A companion
development~\cite{ErikssonObridge} proves the operator-theoretic criterion that
would consume such a transfer operator once it exists; the two are not yet
joined, and we do not claim they are.

\section{The gauge system}

Let $N \ge 1$ and let $\mathrm{Edge}$, $\mathrm{Plaq}$ be finite types.  A
\emph{gauge system} is an incidence $\mathrm{inc} : \mathrm{Plaq} \to
\mathrm{Edge} \to \Z_N$, a strictly positive single-plaquette Boltzmann factor
$\mathrm{weight} : \Z_N \to \R$, and an involution $\theta$ of the edge set.  A
configuration is $U : \mathrm{Edge} \to \Z_N$; the holonomy of a plaquette is
the incidence-weighted sum
\[
  \mathrm{hol}(U,p) \;=\; \sum_{e} \mathrm{inc}(p,e)\,U(e),
\]
which for an abelian group is the whole content of ``the product around the
boundary''; and the Gibbs weight is
$\mathrm{boltz}(U) = \prod_p \mathrm{weight}(\mathrm{hol}(U,p))$.  Every
expectation is the finite sum
\[
  \mathbb{E}[F] \;=\; \Bigl(\sum_U \mathrm{boltz}(U)\,F(U)\Bigr)\Big/ Z,
  \qquad Z = \sum_U \mathrm{boltz}(U) > 0 ,
\]
so no measure theory occurs anywhere in the development.

Carrying the lattice geometry as data rather than fixing a box means the
analytic content is proved once, for every lattice and every reflection at
once.  The price is that reflection invariance of the weight is a hypothesis
rather than a consequence; it is a property of the geometry, and each concrete
system discharges it.

\section{The analytic step}

\begin{definition}
For a kernel $K : X \times X \to \C$ on a finite type write
\[
  Q_K(F) \;=\; \sum_{x}\sum_{y} \overline{F(x)}\,K(x,y)\,F(y),
\]
and call $K$ \emph{positive semidefinite} when $Q_K(F)$ is a non-negative real
for every $F$.  (Stated in that form, reality is part of the conclusion rather
than a side condition.)
\end{definition}

\begin{definition}[non-negative character combination]
$K$ is a non-negative combination for the family $\varphi$ when $c_i \ge 0$ for
all $i$ and $K(x,y) = \sum_i c_i\,\varphi_i(x)\overline{\varphi_i(y)}$.
\end{definition}

\begin{theorem}
\label{thm:psd}
A non-negative character combination is positive semidefinite.
\end{theorem}

\begin{proof}
Put $w_i = \sum_y \overline{\varphi_i(y)}F(y)$.  Resumming,
$Q_K(F) = \sum_i c_i\,\overline{w_i}\,w_i = \sum_i c_i |w_i|^2 \ge 0$.
\end{proof}

\begin{proposition}
\label{prop:closure}
The class is closed under entrywise products and under conjugation by a
diagonal.
\end{proposition}

\begin{proof}
For products, index by pairs: $\varphi_i\psi_j$ is again a family and
$c_id_j \ge 0$.  For the diagonal, replace $\varphi_i$ by $D\varphi_i$; the
coefficients are unchanged.
\end{proof}

\begin{remark}
Proposition~\ref{prop:closure} is what would otherwise be the Schur product
theorem, and Theorem~\ref{thm:psd} is what would otherwise be one direction of
Bochner's theorem on a finite abelian group.  Defining the class by the shape
of the combination rather than by a condition on Fourier coefficients makes
both elementary, and the development consequently uses no spectral and no
matrix-positivity API at all.
\end{remark}

\section{The geometric step}

\begin{definition}[reflection splitting]
A \emph{splitting} of a gauge system consists of a finite type
$\mathrm{Half}$, an equivalence $\mathrm{Config} \simeq \mathrm{Half} \times
\mathrm{Half}$ under which the reflection acts as the swap, a strictly positive
half-space factor $w$, and a crossing kernel $K$, such that
\[
  \mathrm{boltz}(U) \;=\; w(x)\,w(y)\,K(x,y),
  \qquad (x,y) = \text{the two halves of } U .
\]
\end{definition}

\begin{definition}
For $F$ a function of one half-space configuration, the
Osterwalder--Seiler pairing is
\[
  \langle F \rangle_{\mathrm{OS}} \;=\;
  \sum_U \mathrm{boltz}(U)\,\overline{F(x_U)}\,F(y_U).
\]
\end{definition}

\begin{theorem}[the identity that does the work]
\label{thm:identity}
$\langle F\rangle_{\mathrm{OS}} = Q_{M}(F)$ with $M(x,y) = w(x)K(x,y)w(y)$.
\end{theorem}

\begin{proof}
Reindex the sum over configurations by the splitting and factor the weight.
\end{proof}

\begin{corollary}[reflection positivity]
\label{cor:rp}
If the crossing kernel is a non-negative character combination then
$\langle F\rangle_{\mathrm{OS}}$ is a non-negative real for every $F$.
\end{corollary}

\begin{proof}
By Proposition~\ref{prop:closure} the half-space factors are absorbed into the
combination, and Theorem~\ref{thm:psd} applies to $M$; conclude with
Theorem~\ref{thm:identity}.
\end{proof}

\section{The physical step: where a Boltzmann factor is used}

Sections 3 and 4 assume nothing about any particular weight.  The single place
a property of the Boltzmann factor is used is here.

\begin{lemma}[weights to kernels]
\label{lem:conv}
Let $X$ be a finite abelian group and let $\chi$ be a family with
$\chi_i(x-y) = \chi_i(x)\overline{\chi_i(y)}$.  If $k(u) = \sum_i c_i\chi_i(u)$
with $c_i \ge 0$, then $(x,y)\mapsto k(x-y)$ is a non-negative character
combination.
\end{lemma}

\begin{theorem}[the $\Z_2$ Wilson instance]
\label{thm:z2}
For $\Z_2$ the Wilson factor is $e^{\beta s}$ with $s = \pm 1$.  Its expansion
in the two characters of $\Z_2$ is
\[
  e^{\beta s} \;=\; c_0 + c_1\,\chi_1,
  \qquad
  c_0 = \frac{e^{\beta}+e^{-\beta}}{2},
  \quad
  c_1 = \frac{e^{\beta}-e^{-\beta}}{2},
\]
and $c_0 \ge 0$ always while $c_1 \ge 0$ exactly when $\beta \ge 0$.  Hence for
$\beta \ge 0$ the $\Z_2$ Wilson crossing kernel is a non-negative character
combination, and by Corollary~\ref{cor:rp} the pairing is positive.  For
$\beta > 0$ the coefficient $c_1$ is strictly positive, so the weight is
genuinely non-constant and the instance is not the degenerate $\beta = 0$ one.
\end{theorem}

The coefficients are written in that explicit form, rather than as $\cosh$ and
$\sinh$, so that the mechanized proof needs no hyperbolic-function lemma: the
non-negativity of $c_1$ is monotonicity of $\exp$.

\begin{remark}[the limit of the instance]
\label{rem:limit}
Theorem~\ref{thm:z2} settles $\Z_2$ at non-negative coupling and nothing more.
For $\Z_N$ with $N > 2$ the character coefficients of the Wilson factor are
discrete Bessel-type sums; their non-negativity is a genuine analytic fact and
is \emph{not} established here.  We state it as the natural next target rather
than as folklore we are entitled to.
\end{remark}

\section{Non-vacuity}

The predicates above are of a shape that admits degenerate inhabitants —
positive semidefiniteness is satisfied by the zero kernel, a splitting by a
one-point system — so each is exhibited non-degenerately.  The kernel class is
inhabited by a kernel that is not identically zero.  The splitting structure is
discharged concretely for a system with two edges exchanged by the reflection,
and the whole chain fires on it: reflection positivity is obtained by
\emph{discharging} a splitting and a character combination, not by assuming a
pairing.  The physics content of that particular witness is nil, its weight
being constant, and the module says so; the content is that the machinery
composes.  Theorem~\ref{thm:z2} supplies the non-degenerate weight separately.

\section{What remains}

Within the Osterwalder--Seiler chain the remaining obligations are the GNS
quotient of the pairing proved positive here, the transfer operator with
$0 \le T \le 1$ fixing the vacuum, and the identification of Euclidean
correlators with its matrix elements.  The first of these is within reach of
the library as it stands: a positive semidefinite sesquilinear form is a
\lean{PreInnerProductSpace.Core}, which requires no definiteness, and the
quotient by its null space is the separation quotient of the induced seminorm,
which carries an inner-product-space instance already.  We record that as a
measured observation about the library, not as work done.

We do not claim that list is exhaustive for any particular notion of mass gap,
and we do not claim any of it here.

\section{Reproducibility and verification}

All statements are machine-checked in Lean~4 (toolchain
\lean{leanprover/lean4:v4.29.0-rc6}) against Mathlib pinned to commit
\lean{07642720480157414db592fa85b626dafb71355b}.  The freeze is the source tree
at commit \lean{\anchor}.  Build and oracle figures are recorded in the
repository's verification ledger at that commit.  Every declaration reachable
from the core root is subject to an axiom oracle; the transcript records that
the axiom sets occurring are subsets of
\lean{[propext, Classical.choice, Quot.sound]}, with zero occurrences of
\lean{sorryAx} and no project axioms.

\small
\begin{longtable}{@{}>{\raggedright\arraybackslash}p{0.34\textwidth} >{\raggedright\arraybackslash}p{0.48\textwidth} l@{}}
\toprule
Statement here & Lean name & Mod. \\
\midrule
\endhead
Gauge system, weight, expectation & \lean{GaugeData}, \lean{boltz},
  \lean{expect} & Z \\
Reflection on configurations & \lean{reflConfig}, \lean{ReflInvariant} & Z \\
$Q_K$, positive semidefiniteness & \lean{kernelForm}, \lean{IsPSDKernel} & P \\
Character combination & \lean{IsCharCombo} & P \\
Theorem~\ref{thm:psd} & \lean{isPSDKernel\_of\_charCombo} & P \\
Prop.~\ref{prop:closure} & \lean{isCharCombo\_mul},
  \lean{isCharCombo\_diagConj} & P \\
Splitting & \lean{ReflectionSplitting} & R \\
OS pairing & \lean{osPairing}, \lean{osKernel} & R \\
Theorem~\ref{thm:identity} & \lean{osPairing\_eq\_kernelForm} & R \\
Cor.~\ref{cor:rp} & \lean{osPairing\_nonneg} & R \\
Lemma~\ref{lem:conv} & \lean{isCharCombo\_of\_convolution} & W \\
Theorem~\ref{thm:z2} & \lean{z2Wilson\_isCharCombo},
  \lean{z2Wilson\_isPSDKernel}, \lean{z2Coeff\_one\_pos} & W \\
Non-vacuity of the chain & \lean{minimalSplitting\_reflectionPositive} & R \\
\bottomrule
\end{longtable}
\normalsize

Modules: Z = \osfile{ZNSubstrate}{YangMills/OS/ZNSubstrate.lean};
P = \osfile{PSDKernel}{YangMills/OS/PSDKernel.lean};
R = \osfile{ReflectionSplitting}{YangMills/OS/ReflectionSplitting.lean};
W = \osfile{WilsonCharCombo}{YangMills/OS/WilsonCharCombo.lean}.
Core root: \repofile{YangMillsCore.lean}{YangMillsCore.lean}; oracle driver:
\repofile{oracle\_check.lean}{oracle\_check.lean}; governance and scope record:
\repofile{docs/O-BRIDGE-CHARTER.md}{docs/O-BRIDGE-CHARTER.md}.

\subsection*{Provenance}

Reflection positivity for lattice gauge theories is due to Osterwalder and
Seiler~\cite{OsterwalderSeiler}, and the character-expansion criterion for
abelian models is classical; \textbf{no mathematical novelty is claimed for
either}.  What is offered is the mechanization, and the separation of the
argument into an analytic step, a geometric step, and a single point of
contact with a Boltzmann factor, which is what makes the $\Z_N$, $N>2$ gap in
Remark~\ref{rem:limit} visible as a precise open statement rather than a
hand-wave.

\begin{thebibliography}{9}

\bibitem{OsterwalderSeiler}
K.~Osterwalder and E.~Seiler,
\emph{Gauge field theories on a lattice},
Ann. Physics \textbf{110} (1978), 440--471.

\bibitem{GlimmJaffe}
J.~Glimm and A.~Jaffe,
\emph{Quantum Physics: A Functional Integral Point of View},
2nd ed., Springer, 1987.

\bibitem{ErikssonObridge}
L.~Eriksson,
\emph{Clustering and the transfer-operator gap: a machine-checked
dense-family criterion}, 2026.

\bibitem{Lean4}
L.~de~Moura and S.~Ullrich,
\emph{The Lean 4 theorem prover and programming language},
CADE~28 (2021), 625--635.

\bibitem{Mathlib}
The mathlib Community,
\emph{The Lean mathematical library},
CPP 2020, 367--381.

\end{thebibliography}

\end{document}
